\SourceLeaf{360}%
verzichtet also auf die Vorteile, die sich aus der kleineren Entfernung, dem in der Praxis ungleich häufigeren Fall, mit anderen Methoden ziehen lassen. Die \textsc{Bessel}sche Methode wird sich daher sinnvoll nur anwenden lassen, wenn die beiden Punkte über eine längere Kette von Dreiecken miteinander verbunden sind und nicht unmittelbar durch die Seite eines tatsächlich gemessenen Dreiecks. Dann muss man aber mit Recht fragen. wie man die \textit{Daten} für die Aufgabe erhalten soll, nämlich die wirkliche Länge der geodätischen Linie zwischen beiden Punkten und den Winkel, den sie an einem Endpunkt mit dem Meridian bildet? Diese Bestimmung durch eine bloß \textit{sphärische} Rechnung der Übergangsdreiecke vorzunehmen, wie \textsc{Bessel} es bei der wenig ausgedehnten preußischen Gradmessung getan hat, wäre bei viel größerer Entfernung nicht mehr zulässig. Soll dieser Übergang aber sphäroidisch berechnet werden, so erfordert das allein schon ebenso viel Arbeit wie die unmittelbare Bestimmung der Breite, der Länge und des Rückazimuts für jeden folgenden Punkt. Diese Bemerkungen gelten auch für die Lösungsmethode von \textsc{Ivory}, die der von \textsc{Bessel} sehr ähnlich ist, den eigentlichen praktischen Bedarf aber wenig berücksichtigt.

Über die Methode der vorliegenden Abhandlung sei hier noch Folgendes bemerkt.

Die Formeln geben unmittelbar die \textit{Differenzen} zwischen den beiden Breiten und zwischen den beiden Azimuten sowie den Längenunterschied an. Auf der Kleinheit dieser Differenzen beruht die Genauigkeit der Rechnung ohne Logarithmen mit mehr Stellen, sofern bei den Azimuten das eine von der Südseite und das andere von der Nordseite des Meridians aus gezählt wird. Die Symmetrie und Einfachheit der Formeln beruht dagegen darauf, dass sie zunächst vom Mittel der beiden Breiten und vom Mittel der beiden Azimute abhängen, statt von der Breite und dem Azimut an einem Endpunkt. Daraus folgt, dass sich die Formeln zur Lösung der oben formulierten Aufgabe nur durch ein indirektes Verfahren oder durch schrittweise Annäherung anwenden lassen. Geübte Rechner, die mit den Hilfsmitteln solcher Rechenschritte vertraut sind, werden darin kaum eine Schwierigkeit sehen, zumal fast immer angenommen werden kann, dass bei der genauen %
